
\def\PUDcodigo{ACE007}

\if\COMPILAR0
   \def\PUDcodacad{04507.33}
\else
   \def\PUDcodacad{04506.32}
\fi

\def\PUDdisciplina{Ética e Cidadania}

\def\PUDsemestre{S7}

\def\PUDhoras{40}

%NOVOS ---------------------------------------------------
\def\PUDhorasTeoria{40}
\def\PUDhorasPratica{0}

\if\COMPILAR0
   \def\PUDinterECA{\hyperref[PUD:ACE017]{Introdução à Engenharia (04507.5)}}
\else
   \def\PUDinterECA{\hyperref[PUD:ACE017]{Introdução à Engenharia (04506.5)}}
\fi

\def\PUDatuaisECA{---}

\def\PUDrecursos{Projetor multimídia; Quadro branco e pincel.}

%----------------------------------------------------------

\def\PUDcreditos{2}

\def\PUDprerequisito{---}

\def\PUDpreacad{---}

\def\PUDobjetivos{Familiarizar-se com os princípios da ética profissional no âmbito das organizações e sua importância para a transformação da sociedade; Conhecer formas de análise e aplicação dos códigos de ética profissionais, com ênfase no do engenheiro. Familiarizar-se com os princípios da ética ambiental.}

\def\PUDementa{Responsabilidade social do engenheiro, Profissão do homem diante da participação, Código de Ética Profissional, Os órgãos de representação de classe, Princípios gerais de legislação trabalhista, Direito sindical, Seguridade social. Ética ambiental. }

\def\PUDconteudo{ &  A responsabilidade social do engenheiro: Valor social da profissão;  Responsabilidade social da profissão;  Função social do engenheiro;  Deveres profissionais;  Atualização constante e aperfeiçoamento cultural;  Influência das realizações profissionais no ambiente e na sociedade; 
\\ 
 &  Profissão do homem diante da participação: Participação do engenheiro na comunidade local, nacional ou internacional;  Relação do engenheiro com outros profissionais;  
\\ 
 &  Código de Ética Profissional: Elemento de Ética;  Base filosófica do Código de Ética Profissional;  Atitude profissional;  Virtudes básicas;  Virtudes específicas da profissão;  Código de Ética Profissional do engenheiro;  Julgamento da conduta ética na classe.  
\\ 
 &  Órgãos de classe: CONFEA, CREA e Câmaras Especializadas;  Outros órgãos de classe;  Lei de regulamentação da profissão do engenheiro;  Anotação de Responsabilidade Técnica - ART;  
\\ 
 &  Noções de legislação trabalhista.  
\\ 
 &  Noções de direito sindical.  
\\ 
 &  Noções de seguridade social.
 \\
  & Ética ambiental no âmbito da engenharia. }

\def\PUDmetodo{Aulas expositórias que podem ser teóricas e/ou práticas, onde as práticas no laboratório serão marcadas no decorrer da disciplina.}

\def\PUDavalia{A avaliação será feita com aplicação de provas teóricas e/ou práticas, além da possibilidade de inclusão de trabalhos e seminários no decorrer da disciplina.}

\def\PUDbibbasica{ &  \href{www.biblioteca.ifce.edu.br/index.asp?codigo_sophia=61932}{Srour}, Robert Henry.   Ética Empresarial.  4o ed.  Rio de Janeiro, RJ: Elsevier, 2013.  213p.  ISBN: 9788535264470.
\\ 
 &   \href{www.biblioteca.ifce.edu.br/index.asp?codigo_sophia=62399}{Sá}, Antônio L.  Ética profissional.  9o ed.  São Paulo, SP: Atlas, 2015.  312p.  ISBN: 9788522455348.
\\ 
 &   \href{www.biblioteca.ifce.edu.br/index.asp?codigo_sophia=24317}{Dias}, Reinaldo.  Marketing ambiental: ética, responsabilidade social e competitividade nos negócios.  São Paulo, SP: Atlas, 2009.  200p.  ISBN: 9788522446766.  
 %&   SÁNCHEZ VÁZQUEZ, Adolfo.   Ética.  18ª edição.  Ed.  Rio de Janeiro: Civilização Brasileira, 1998. 
 }

\def\PUDbibcomplementar{ & \href{www.biblioteca.ifce.edu.br/index.asp?codigo_sophia=39514}{Holtzapple}, Mark Thomas.  Introdução à engenharia. Rio de Janeiro, RJ : LTC, 2013.  ISBN: 9788521615118.
%  SINGER, Peter.  Ética Prática.  São Paulo: Martins Fontes, 1994.  ISBN 9788533616684
\\ 
 &  \href{www.biblioteca.ifce.edu.br/index.asp?codigo_sophia=40790}{Barros}, Benjamim F. de.  NR-10: guia prático de análise e aplicação.  3º ed.  São Paulo, SP : Érica, 2014. 204p.  ISBN: 9788536502748.
%    BENNETT, Willian J.  O livro das Virtudes II.  Rio de janeiro.  Ed.  Nova Fronteira, 1996. 
\\ 
 &  \href{www.biblioteca.ifce.edu.br/index.asp?codigo_sophia=24014}{Pepplow}, Luiz Amilton.  Segurança do trabalho.  Curitiba, PR : Base Editorial, 2010. 256p. ISBN: 9788579055430.
\\
 &  \href{www.biblioteca.ifce.edu.br/index.asp?codigo_sophia=56117}{Gallo}, Sílvio.  Ética e cidadania: Caminhos da filosofia.  E-book.  Papirus.  116p.  ISBN: 9788530811525.
\\
 &  \href{www.biblioteca.ifce.edu.br/index.asp?codigo_sophia=58459}{Alencastro}, Mario Sergio Cunha.  Ética e meio ambiente: construindo as bases para um futuro sustentável.  E-book.  Intersaberes.  186p.  ISBN: 9788544301173. }
%   Código de Ética Profissional do Servidor Público Civil do Poder Executivo Federal, Decreto N. 1.171 de 22 de junho de 1994, fonte: \url{http://www.planalto.gov.br/ccivil_03/decreto/d1171.htm}. Data: 13/02/2017.
%\url{http://www. ifce. edu. br/images/stories/menu_superior/Ensino/ROD/ROD-Comisso_de_Sistematizao27. pdf

\def\PUDautor{Fabrício Bandeira da Silva}

\def\PUDdata{2014-05-19}

\def\PUDrevisao{2}

\def\PUDrevisor{Samuel Vieira Dias}

\def\PUDdatarevisao{2017-05-23}

\PUDcorpo
